\section{The Bailout Protocol}
\label{sec:bailout-protocol}

The Bailout Protocol is the runtime mechanism by which the BX3 Framework enforces its upstream accountability guarantee: that every unresolvable exception state propagates to a human accountability anchor, bypassing all intermediate machine actors, with complete diagnostic context, and that the system never fails downward into algorithmic chaos. The protocol is triggered not by machine error but by the Bounds Engine's correct identification of the boundary of its own authority. It is an \textit{authority-exceeded} signal, not a failure signal.

This section specifies the protocol in full formal detail: its state machine, its trigger conditions, its escalation algorithm, its bypass mechanism, and its timeout handling. The specification is complete enough to serve as an authoritative reference for implementation, audit, and formal verification.

\subsection{Formal Definitions}
\label{sec:bailout-formal-definitions}

We define the following entities and invariants that govern the Bailout Protocol:

\medskip

\noindent\textbf{Definition 1 (Bounds Engine Node).} A Bounds Engine node $B$ is a tuple:
\begin{equation}
B \triangleq \langle C_B, M_B, P_B, \text{BailoutSM}_B \rangle
\end{equation}
where:
\begin{itemize}\setlength\itemsep{0.25em}
\item $C_B$ is the constraint set governing $B$'s operational scope, authored by a Purpose Layer owner.
\item $M_B \subseteq C_B$ is the active mandate under the present operational context.
\item $P_B$ is the parent node in the system hierarchy; if $B$ is a root node, $P_B = \texttt{null}$.
\item $\text{BailoutSM}_B$ is the per-instance Bailout Protocol state machine for node $B$, with state $\sigma_B \in \Sigma$, where $\Sigma = \{\text{IDLE}, \text{BOUNDED}, \text{BAIL\_PENDING}, \text{ESCALATING}, \text{HUMAN\_ACKNOWLEDGED}, \text{RESOLVED}\}$.
\end{itemize}

\medskip

\noindent\textbf{Definition 2 (Constraint Set).} The constraint set $C_B$ is a set of predicate functions over the system state $S$:
\begin{equation}
C_B = \{ c_i : S \rightarrow \{\texttt{true}, \texttt{false}\} \mid i = 1, \ldots, n \}
\end{equation}
A proposed action $a$ is \textit{approved} by $B$ if $\forall c_i \in M_B : c_i(S \mid a) = \texttt{true}$. A proposed action $a$ is \textit{rejected} by $B$ if $\exists c_i \in M_B : c_i(S \mid a) = \texttt{false}$.

\medskip

\noindent\textbf{Definition 3 (Unresolvable State).} A system state $s \in S$ is \textit{unresolvable} by $B$ under mandate $M_B$ if and only if:
\begin{equation}
\text{isUnresolvable}(s, M_B) \triangleq \nexists a : \forall c_i \in M_B : c_i(s \mid a) = \texttt{true}
\end{equation}
That is, no proposed action exists that satisfies all active constraints. This is the foundational condition for triggering the Bailout Protocol.

\medskip

\noindent\textbf{Definition 4 (BailoutTrigger).} When a trigger condition is satisfied at node $B$, the Bounds Engine constructs a BailoutTrigger object $T$:
\begin{equation}
T \triangleq \langle \text{id}, B, s, M_B, \kappa, \psi, \tau, \phi \rangle
\end{equation}
where:
\begin{itemize}\setlength\itemsep{0.25em}
\item $\text{id}$: unique identifier for this bailout event.
\item $B$: the originating Bounds Engine node.
\item $s$: the system state at time of triggering.
\item $M_B$: the active mandate at time of triggering.
\item $\kappa \in \mathbb{K}$: the trigger condition class (Section~\ref{sec:bailout-trigger-conditions}).
\item $\psi$: the diagnostic snapshot of $B$'s internal state at time of firing.
\item $\tau \in \mathbb{R}^+$: the Unix timestamp at which the trigger was fired.
\item $\phi = [\,]$: the propagation chain, initially empty; appended to as $T$ propagates up the hierarchy.
\end{itemize}

\medskip

\noindent\textbf{Definition 5 (Trigger Condition Class).} The set $\mathbb{K}$ of recognized trigger condition classes is:
\begin{equation}
\mathbb{K} = \{
\text{CONSTRAINT\_UNCOVERED},
\text{UNRESOLVABLE\_STATE},
\text{INTERCONSTRAINT\_CONFLICT},
\text{FACT\_LAYER\_VIOLATION},
\text{SCOPE\_EXCEEDED},
\text{MANDATE\_MODIFICATION\_REQUEST}
\}
\end{equation}
Each class corresponds to a distinct boundary condition of authority (Section~\ref{sec:bailout-trigger-conditions}).

\medskip

\noindent\textbf{Definition 6 (Purpose Layer Owner).} The Purpose Layer owner of node $B$, denoted $\text{owner}(B)$, is the named human actor who authored $B$'s mandate $M_B$ and bears accountability for $B$'s outcomes. $\text{owner}(B)$ is the sole permissible terminal recipient of a BailoutTrigger.

\medskip

\noindent\textbf{Definition 7 (Hard Block).} When $\sigma_B \in \{\text{BAIL\_PENDING}, \text{ESCALATING}, \text{HUMAN\_ACKNOWLEDGED}\}$, the Fact Layer enforces a hard block $\Gamma_B$ on all physical actuator commands originating from $B$. Formally:
\begin{equation}
\forall \text{actuator } \alpha : \forall \text{command } a \in A_\alpha \text{ proposed by } B : \Gamma_B(\sigma_B) = \texttt{block}(a) \quad \text{if } \sigma_B \neq \text{IDLE} \land \sigma_B \neq \text{RESOLVED}
\end{equation}
The hard block is monotonic: applied at BOUNDED, not released until RESOLVED. No confidence signal, timer, or business context can override $\Gamma_B$.

\medskip

\noindent\textbf{Definition 8 (Timeout Constants).} The Bailout Protocol defines two timeout constants enforced by the Fact Layer:
\begin{equation}
T_{\text{bound}} = 300 \quad \text{(seconds, 5 minutes)} \qquad T_{\text{human}} = 7200 \quad \text{(seconds, 2 hours)}
\end{equation}
$T_{\text{bound}}$ is the maximum duration a trigger may remain in the BOUNDED state before automatic transition to BAIL\_PENDING. $T_{\text{human}}$ is the interval at which the protocol issues a staleness warning after entering HUMAN\_ACKNOWLEDGED.

\subsection{State Machine: Formal Specification}
\label{sec:bailout-state-machine}

The Bailout Protocol operates as a deterministic finite state machine $\text{BailoutSM}_B$, hard-encoded in the Fact Layer for each Bounds Engine node $B$. The state machine has six mutually exclusive, collectively exhaustive states. The state $\sigma_B$ is stored in the Fact Layer worksheet and is the authoritative source of truth for the node's current protocol phase.

\bigskip

\begin{table}[htbp]
\centering
\caption{Bailout Protocol State Machine $\mathcal{B}$: States and Invariants}
\label{tab:bailout-states}
\begin{tabular}{p{0.18\linewidth} | p{0.38\linewidth} | p{0.36\linewidth}}
\hline
\textbf{State} & \textbf{Description} & \textbf{State Invariant} \\
\hline
$\sigma = \text{IDLE}$ & Default operational state. No trigger is active. Bounds Engine evaluates proposed actions within mandate $M_B$. & $\Gamma_B = \texttt.pass}$. No hard block active. No $T$ exists. \\
\hline
$\sigma = \text{BOUNDED}$ & Trigger condition detected. BailoutTrigger object $T$ constructed. Diagnostic context $\psi$ being built. Hard block applied. & $\Gamma_B = \texttt{block}$. No actuator commands forwarded. Timer $t_{\text{bound}}$ running. \\
\hline
$\sigma = \text{BAIL\_PENDING}$ & Diagnostic context complete. $T$ sealed and ready for propagation. Hard block maintained. & $\Gamma_B = \texttt{block}$. $T$ is immutable. No machine actor may resolve $T$. Next transition must be to ESCALATING. \\
\hline
$\sigma = \text{ESCALATING}$ & $T$ dispatched to parent node $P_B$. Trigger in transit. Hard block maintained on originating node. & $\Gamma_B = \texttt{block}$. Propagation chain $\phi$ updated atomically. Machine nodes must forward without attempting resolution. \\
\hline
$\sigma = \text{HUMAN\_ACKNOWLEDGED}$ & $T$ received by $\text{owner}(B)$. Human anchor actively reviewing. Hard block maintained. & $\Gamma_B = \texttt{block}$. Timer $t_{\text{human}}$ running. Only $\text{owner}(B)$ may issue resolution. \\
\hline
$\sigma = \text{RESOLVED}$ & Purpose Layer owner issued resolution. Forensic Ledger entry written. Hard block released. State machine for this $T$ terminates. & $\Gamma_B = \texttt.pass}$. Resolution recorded. No outbound transitions. \\
\hline
\end{tabular}
\end{table}

\bigskip

The state machine transition function is:
\begin{equation}
\delta : \Sigma \times E \rightarrow \Sigma
\end{equation}
where $E$ is the event set. All transitions are deterministic events enforced by the Fact Layer pre-commit hook. Any event not in the defined transition set is rejected.

\subsection{Transition Conditions}
\label{sec:bailout-transitions}

Transitions are deterministic events --- not timers, confidence thresholds, or probabilistic triggers. The complete transition function is defined as follows:

\bigskip

\noindent\textbf{From IDLE:}
\begin{equation}
\delta(\text{IDLE}, \texttt{TRIGGER\_CONDITION\_SATISFIED}) = \text{BOUNDED}
\end{equation}
Fires when the Bounds Engine detects that a trigger condition $\kappa \in \mathbb{K}$ is satisfied at node $B$. This is the sole transition out of IDLE and the entry point for every Bailout Protocol instance.

\bigskip

\noindent\textbf{From BOUNDED:}
\begin{equation}
\delta(\text{BOUNDED}, \texttt{DIAGNOSTIC\_COMPLETE}) = \text{BAIL\_PENDING}
\end{equation}
Fires when the Bounds Engine has completed construction of the full BailoutTrigger object $T$, including the diagnostic snapshot $\psi$. The transition is atomic: the Fact Layer commits both the state change and the sealed $T$ simultaneously.
\begin{equation}
\delta(\text{BOUNDED}, \texttt{TIMER\_EXPIRED\_BOUND}) = \text{BAIL\_PENDING}
\end{equation}
Fires when elapsed time $t_{\text{bound}}$ since entering BOUNDED exceeds $T_{\text{bound}}$. The Fact Layer forces this transition, sets \texttt{CONSTRUCTION\_TIMEOUT} in $T$, and logs the timeout event in the Forensic Ledger. $T$ proceeds to escalation with whatever diagnostic context was available at timeout.

\bigskip

\noindent\textbf{From BAIL\_PENDING:}
\begin{equation}
\delta(\text{BAIL\_PENDING}, \texttt{TRIGGER\_DISPATCHED}) = \text{ESCALATING}
\end{equation}
Fires immediately upon dispatch of $T$ to the parent node $P_B$. The propagation chain $\phi$ is updated atomically with the dispatch event. No machine actor may delay, absorb, or override this transition.

\bigskip

\noindent\textbf{From ESCALATING:}
\begin{equation}
\delta(\text{ESCALATING}, \texttt{RECEIVED\_BY\_HUMAN\_ANCHOR}) = \text{HUMAN\_ACKNOWLEDGED}
\end{equation}
Fires when the parent node $P_B$ is confirmed to be $\text{owner}(B)$ (a Purpose Layer human anchor). This is the terminal propagation transition.
\begin{equation}
\delta(\text{ESCALATING}, \texttt{RECEIVED\_BY\_MACHINE\_NODE}) = \text{ESCALATING}
\end{equation}
Fires when $P_B$ is another Bounds Engine node. $T$ is re-dispatched to $P_{P_B}$, and $\phi$ is appended. This transition repeats recursively until a human anchor is reached.

\bigskip

\noindent\textbf{From HUMAN\_ACKNOWLEDGED:}
\begin{equation}
\delta(\text{HUMAN\_ACKNOWLEDGED}, \texttt{HUMAN\_RESOLUTION\_ISSUED}) = \text{RESOLVED}
\end{equation}
Fires when $\text{owner}(B)$ issues a resolution. The Fact Layer records the resolution in the Forensic Ledger, releases $\Gamma_B$, and transitions to IDLE. This is the sole terminal transition.
\begin{equation}
\delta(\text{HUMAN\_ACKNOWLEDGED}, \texttt{TIMER\_EXPIRED\_HUMAN}) = \text{HUMAN\_ACKNOWLEDGED}
\end{equation}
Fires when elapsed time $t_{\text{human}}$ since entering HUMAN\_ACKNOWLEDGED exceeds $T_{\text{human}}$. The state does not change; this transition re-fires at intervals of $T_{\text{human}}$ until a resolution is issued.

\bigskip

\noindent\textbf{From RESOLVED:} There are no outbound transitions. The state machine for this trigger instance terminates.

\bigskip

The Fact Layer enforces that no transition not explicitly defined above may occur. Any attempt to resolve a trigger at an unauthorized state (e.g., $\delta(\text{BOUNDED}, \texttt{MACHINE\_RESOLUTION\_ATTEMPTED})$) is undefined and therefore impermissible: the Fact Layer rejects it and logs the attempt as a \texttt{BYPASS\_VIOLATION} in the Forensic Ledger.

\subsection{Bail-Out Trigger Conditions}
\label{sec:bailout-trigger-conditions}

The Bailout Protocol is triggered by the Bounds Engine when it correctly identifies a condition that exceeds its resolution authority. The trigger is an \textit{authority boundary} signal, not a failure signal. A functioning Bounds Engine fires bailouts precisely because it has correctly identified an out-of-scope condition.

\bigskip

A Bounds Engine node $B$ fires a BailoutTrigger when any of the following conditions is satisfied:

\bigskip

\noindent\textbf{T1 (Constraint Uncovered).} $B$ receives a request to perform action $a$, and $a$ is not governed by any constraint in $M_B$:
\begin{equation}
\nexists c_i \in M_B : c_i \text{ governs } a
\end{equation}
This condition indicates that the proposed action is outside the scope of the current mandate and requires Purpose Layer authorization.

\bigskip

\noindent\textbf{T2 (Unresolvable State).} $B$ detects a system state $s$ for which $\text{isUnresolvable}(s, M_B) = \texttt{true}$ (Definition 3). The environment has moved into a configuration that the current mandate has no approved response to --- not because the Bounds Engine failed, but because the mandate has a genuine gap.

\bigskip

\noindent\textbf{T3 (Interconstraint Conflict).} $B$ detects a conflict between two or more active constraints $c_i, c_j \in M_B$ such that:
\begin{equation}
\nexists a : c_i(s \mid a) = \texttt{true} \land c_j(s \mid a) = \texttt{true}
\end{equation}
The mandate contains internally contradictory requirements and requires Purpose Layer disambiguation or mandate revision.

\bigskip

\noindent\textbf{T4 (Fact Layer Violation).} $B$ receives a proposed action $a$ that, when validated by the Fact Layer enforcement model, violates a Fact Layer enforcement rule $r \in R_{\text{fact}}$:
\begin{equation}
\text{FactLayer\_validate}(a, r) = \texttt{violation}
\end{equation}
The Bounds Engine's proposal would breach a deterministic safety constraint that the Fact Layer enforces independently.

\bigskip

\noindent\textbf{T5 (Scope Exceeded).} $B$ receives a request to perform action $a$ that falls within its technical capability but outside the scope of $M_B$:
\begin{equation}
a \in \text{capability}(B) \land a \notin \text{authorized\_scope}(M_B)
\end{equation}
This condition distinguishes between what $B$ \textit{can} do and what $B$ \textit{is authorized} to do.

\bigskip

\noindent\textbf{T6 (Mandate Modification Request).} $B$ receives a request to modify $C_B$ or $M_B$, which it cannot authorize unilaterally:
\begin{equation}
\text{request\_type}(r) = \texttt{MANDATE\_MODIFICATION}
\end{equation}
Mandate modifications require Purpose Layer approval and route the request to $\text{owner}(B)$.

\bigskip

When a trigger condition is satisfied, the Bounds Engine transitions to BOUNDED immediately. There is no gating condition, no confidence threshold, and no opportunity for the Bounds Engine to attempt an autonomous resolution first. The guarantee of upstream accountability depends on this immediacy: any delay between condition detection and trigger firing creates a window in which an incorrect autonomous action could be taken.

\subsection{Escalation Algorithm}
\label{sec:bailout-escalation}

The escalation algorithm determines the path a BailoutTrigger takes from the originating Bounds Engine node $B$ to the Purpose Layer human anchor $\text{owner}(B)$. The algorithm is defined as a deterministic traversal over the system hierarchy:

\bigskip

\begin{algorithm}[htbp]
\caption{BailoutTrigger Escalation Algorithm}
\label{alg:bailout-escalation}
\begin{enumerate}\setlength\itemsep{0.2em}\setlength\topsep{0.2em}\setlength\partopsep{0em}
\Function{Escalate}{$T, B$}
  \State $\phi \leftarrow T.\phi$ \Comment{Initialize propagation chain}
  \State $B_{\text{current}} \leftarrow B$ \Comment{Start at originating node}
  \While{$B_{\text{current}} \neq \texttt{null}$}
    \State $\sigma_{B_{\text{current}}} \leftarrow \text{ESCALATING}$
    \State $\phi.\text{append}(\langle B_{\text{current}}, \tau_{\text{now}}, \texttt{DISPATCHED} \rangle)$
    \State $T.\phi \leftarrow \phi$ \Comment{Atomic update of propagation chain}

    \If{$\text{isHumanAnchor}(B_{\text{current}})$}
      \State $\sigma_{B_{\text{current}}} \leftarrow \text{HUMAN\_ACKNOWLEDGED}$
      \State $\phi.\text{append}(\langle B_{\text{current}}, \tau_{\text{now}}, \texttt{HUMAN\_RECEIVED} \rangle)$
      \State $T.\phi \leftarrow \phi$
      \State \Return \texttt{SUCCESS}
    \EndIf

    \State $B_{\text{parent}} \leftarrow P_{B_{\text{current}}}$

    \If{$B_{\text{parent}} = \texttt{null}$}
      \State $\phi.\text{append}(\langle \texttt{null}, \tau_{\text{now}}, \texttt{NO\_ANCHOR\_FOUND} \rangle)$
      \State $T.\phi \leftarrow \phi$
      \State \textbf{raise} \texttt{HIERARCHY\_ERROR}: No human anchor in ancestry chain
      \State \Return \texttt{FAILURE}
    \EndIf

    \State $B_{\text{current}} \leftarrow B_{\text{parent}}$
  \EndWhile

  \State \textbf{raise} \texttt{ESCALATION\_LOOP\_ERROR}: Algorithm did not terminate
  \State \Return \texttt{FAILURE}
\EndFunction
\end{enumerate}
\end{algorithm}

\bigskip

The algorithm has the following critical properties:

\begin{itemize}\setlength\itemsep{0.25em}
\item \textbf{Termination.} The system hierarchy is a finite rooted tree. Since $B_{\text{current}}$ strictly ascends the parent chain on each iteration and tree depth is finite, the while loop terminates in at most $d$ iterations, where $d$ is the depth of $B$ in the hierarchy. Divergence is not possible.

\item \textbf{No intermediate resolution.} The algorithm performs no autonomous reasoning and makes no attempt to resolve the trigger at any machine node. Each machine node in the chain is a transit point, not a decision point. This is the commitment that the bypass mechanism enforces.

\item \textbf{Single destination.} The algorithm routes to exactly one human anchor: $\text{owner}(B)$, the Purpose Layer owner of the originating node. It does not clone the trigger to multiple recipients, does not broadcast to a team, and does not select a recipient based on availability. This ensures a single accountable decision-maker.

\item \textbf{Complete traceability.} The propagation chain $\phi$ records every node $T$ passed through, with a timestamp and action at each hop. This record is included in the Forensic Ledger entry for this bailout event and solves the abstraction leakage problem.
\end{itemize}

The Fact Layer enforces the escalation algorithm's invariants at the architectural level. A machine node cannot refuse to forward a BailoutTrigger, cannot select a different parent node, and cannot substitute a machine fallback for the human anchor. The algorithm is not a policy that may be followed or deviated from; it is a physical constraint of the layer architecture.

\subsection{Bypass Mechanism: Why the Protocol Bypasses All Machine Actors}
\label{sec:bailout-bypass}

The most critical property of the Bailout Protocol is that it bypasses all machine actors in its escalation path. When a BailoutTrigger reaches a machine node during escalation, that node does not attempt to resolve the trigger. It forwards the trigger to its parent and stops. The trigger does not request the machine node's opinion, does not accept a confidence-weighted assessment, and does not allow the machine node to substitute its own resolution.

This bypass is not a convention. It is a structural property enforced by the layer architecture. Three failure modes clarify why the bypass is architecturally necessary:

\bigskip

\noindent\textbf{Failure Mode 1: Autonomous Override.} If a machine node in the escalation path could resolve a BailoutTrigger, it would become a de facto autonomous override of human accountability. The node could apply its own reasoning, propose a resolution, and either dismiss the trigger or substitute its own action --- precisely the pattern the Bailout Protocol is designed to prevent. The upstream accountability guarantee would be voided by a silent machine override the human never sees.

\bigskip

\noindent\textbf{Failure Mode 2: Confidence-Based Routing.} If machine nodes could assess a BailoutTrigger and decide whether to escalate based on their own confidence, the system would introduce a probabilistic component into an architectural safety mechanism. A machine node highly confident in its own resolution would suppress the trigger and proceed autonomously --- the failure mode that formal verification requirements in Simplex-style architectures are designed to prevent. The confidence signal is internal to the machine layers and is not a permitted input to the escalation routing decision.

\bigskip

\noindent\textbf{Failure Mode 3: Abstraction Leakage Amplification.} If machine nodes could respond to a BailoutTrigger, the human accountability anchor would lose additional context with each hop, as each intervening machine node transforms, filters, or summarizes diagnostic information before passing it up. This amplification of abstraction leakage would undermine the forensic standard the protocol is designed to maintain. The bypass ensures that the diagnostic context $\psi$ reaching the human anchor is the original context constructed at the originating node, not a transformed version.

\bigskip

The bypass mechanism is enforced at the state machine level. When $\sigma_B = \text{ESCALATING}$, the Fact Layer rejects any inbound event from a machine actor attempting to resolve the trigger, maintains $\sigma_B$ at ESCALATING, and logs the attempt as a \texttt{BYPASS\_VIOLATION} event. This is the same deterministic enforcement mechanism that enforces all Fact Layer constraints.

The architectural implication is that the Bailout Protocol does not merely \textit{permit} escalation to a human anchor --- it \textit{prevents} all alternative termination points. The system cannot accidentally terminate a bailout at a machine actor, because the Fact Layer will not accept a machine resolution at any state other than RESOLVED, and RESOLVED can only be entered via the HUMAN\_ACKNOWLEDGED $\rightarrow$ RESOLVED transition initiated by the confirmed human anchor. The state machine is a physical enforcement mechanism, not a workflow suggestion.

Accountability for consequential decisions in the real world is a social and legal institution that cannot be instantiated in software. The bypass mechanism ensures that the Bailout Protocol maintains this distinction at runtime, under all conditions including system stress, high workload, and component degradation.

\subsection{Timeout Handling}
\label{sec:bailout-timeouts}

The Bailout Protocol defines two timeout constants, $T_{\text{bound}}$ and $T_{\text{human}}$, that govern the maximum duration a trigger may spend in the BOUNDED and HUMAN\_ACKNOWLEDGED states, respectively. Timeout handling prevents indefinite stalling at either phase while preserving the human's full authority to review at their own pace.

\bigskip

\noindent\textbf{BOUNDED State Timeout ($T_{\text{bound}} = 300$ seconds).} The BOUNDED state is the phase during which the Bounds Engine constructs the complete diagnostic context $\psi$ for the BailoutTrigger. In most cases, construction completes within seconds. However, in complex scenarios involving nested system state, multiple sensor inputs, and cross-plane diagnostic data, construction may take longer.

If $T_{\text{bound}}$ expires before construction is complete, the Fact Layer forces the transition to BAIL\_PENDING with the \texttt{CONSTRUCTION\_TIMEOUT} flag set in $T$. The trigger proceeds to escalation with whatever diagnostic context was available at timeout. The Purpose Layer owner is informed that the diagnostic context is incomplete and may request additional information before issuing a resolution.

The purpose of $T_{\text{bound}}$ is not to penalize the Bounds Engine for slow construction, but to prevent a pathological scenario in which a system under stress (where bailout triggering is most likely) stalls indefinitely in the diagnostic phase due to degraded performance. The timeout ensures that escalation begins within a bounded time regardless of system conditions.

\bigskip

\noindent\textbf{HUMAN\_ACKNOWLEDGED State Timeout ($T_{\text{human}} = 7200$ seconds).} The HUMAN\_ACKNOWLEDGED state is the phase during which the Purpose Layer owner reviews the BailoutTrigger and issues a resolution. Human review is intentionally unbounded in the base protocol: there is no scenario in which a machine actor should force a resolution over the non-action of the human accountability anchor. The $T_{\text{human}}$ timeout exists to address notification failure, absence, or oversight --- not to substitute machine judgment for human judgment.

When $T_{\text{human}}$ expires, the Fact Layer:
\begin{enumerate}\setlength\itemsep{0.2em}
\item Logs a \texttt{STALENESS\_WARNING} in the Forensic Ledger for this bailout event.
\item Issues an SMS reminder to $\text{owner}(B)$ via the INBOX routing system, with a direct link to the bailout context and a summary of the triggering condition.
\item Transitions the state machine to HUMAN\_ACKNOWLEDGED (no state change; only a warning is issued).
\end{enumerate}

The $T_{\text{human}}$ timeout may re-fire at intervals of $T_{\text{human}}$ seconds, with each firing logged as a separate \texttt{STALENESS\_WARNING} entry. The Fact Layer continues to log warnings and issue reminders until the human anchor issues a resolution or the system is formally halted by an authorized administrator.

There is no automatic resolution. Even under repeated staleness warnings, the Fact Layer will not accept a resolution from any actor other than the confirmed human anchor. If the human anchor is permanently unreachable, the escalation algorithm's hierarchy traversal surfaces a \texttt{HIERARCHY\_ERROR} condition (Algorithm~\ref{alg:bailout-escalation}, line 20), which escalates the failure to the next available human anchor in the organizational authority chain. This is the only permitted deviation from the single-anchor rule, triggered only by the absence of a human anchor in the ancestry chain --- not by any machine actor's preference.

\bigskip

The timeout handling design reflects the principle that human review is the mandatory terminal state of the Bailout Protocol and that timeouts exist to ensure the system does not stall --- not to substitute machine judgment for human judgment. The guarantees of the Bailout Protocol hold regardless of timeout behavior, because timeouts never permit autonomous resolution.
